\subsection*{一步一步建立直观}

\paragraph{1. 用放大倍率理解复合}

设输入先经过 $g$，再经过 $f$。一个很小的输入位移先被放大为约 $g'(a)h$，随后又被放大为约 $f'(g(a))$ 倍。右图展示其中一台局部放大器：随着取样窗口缩小，曲线附近越来越像直线，而倍率就是切线斜率。



\paragraph{2. 为什么倍率相乘，却不总是常数}

链式法则把两级倍率相乘，但第二级必须在第一级的输出位置测量。移动右图的观察点，切线斜率随位置改变，便能看到“先代入正确位置”不能省略。对每一级都使用同一个输入位置，是复合求导中常见的错误。



\paragraph{3. 乘积的变化来自两块薄条}

乘积可以看作长方形面积。长度和宽度同时改变时，新增面积包含两块一阶薄条，还有一个二阶小角。除以输入步长后，小角项趋零，两块薄条分别留下 $f'g$ 与 $fg'$。右图仅演示差商极限；长方形分解才是乘积法则的几何理由。
